\section{Sweet Home}
\label{sec:sweetHome}

When the PCs enter \npcTroublemaker{}'s home, read the following:

\quoteBox{
	You enter a stone building, dimly illuminated with a few candles. On the upper floor, in a sickbed lays an old woman. \\ \\
	"\npcTroublemaker{}? Is that you, \npcTroublemaker{}?", the woman asks. "Where have you been, you stupid girl? What took you so storming long?"
}

Laying in the bed is \npcGrandma{} \npcGrandmaName{}, a grandmother of \npcTroublemaker{}. She is seldom happy with anything and she always makes her disappointment vocal. In fact, her loud complaints are dominating each conversation.

\roleplayBox{\npcGrandma{} \npcGrandmaName{}}{
	\roleplayTips
		{Grumpy, ungrateful, unlikable}
		{Complain and insult}
		{Elderly and frail lighteye}
		{Bitter and resentful, \npcGrandma{} remember much better days. Nowadays everything is worse and everyone does everything wrong. \npcGrandma{} cannot tolerate this foolishness, making sure every minor error is loudly complained about. She does not have a time to care about the feelings of others. She is a lighteye (of lower dahn) elder, and she deserves only the best! }
}

\reasonBox{
	In the whole scene, there is little to do for the characters. They can of course join conversation, react to a NPC's behavior, etc, but there is no larger goal. The scene is here to paint the picture of \npcTroublemaker{}'s family and it directs characters to \fullref{sec:peoplesHall}. You can modify the scenes from here as long as you achieve those goals. 
}

\reasonBox{
	The point of this location is to ensure the PCs always have a place in Kholinar they can call home and use for resting. Unlikable NPC of \npcGrandma{} introduces moral choice - will they help her during the troubles ahead? Or will they ignore her? While rescuing good and friendly NPCs is an obvious choice, here the PCs can show their character and personality.
}


\subsection{Parshmen Servants}

For this scene, read the following:

\quoteBox{
	A femalen parshmen approaches \npcGrandma{} with a bowl of steaming meal. She kneels next to a bed and starts gently feeding weak woman. However, after the fist spoon, she starts shouting: \\ \\
	"That's disgusting! What are you giving me, stupid chull? Tallew? I am a storming lighteye!" She smashes the bowl on the floor, then adds: "Clean that up, you idiot! Damnation, you are so useless, I paid too much for you!"
}

\npcSingerWorkerGranM[s]{} and \npcSingerWorkerGranF[s]{} are two parshmen working for \npcGrandma{}. Suffice to say - nobody is happy about this arrangement...

Give the characters a moment to interact with the PCs, perhaps comfort the parshmen, or try to calm \npcGrandma{} down.

\reasonBox{
	\npcGrandma{}'s terrible treatment of her parshmen workers will reap the consequences in \fullref{ssec:enemyRebornServants}.
}


\subsection{Unhospitable Host}

After \npcTroublemaker{} invites the party to eat something, read the following:

\quoteBox{
	As the malen parshmen starts preparing food for the dinner, \npcGrandma{} shouts from her bed: \\ \\
	"Hey, what are you doing? This is our food, paid for with our money! We are not doing beggar's feast in here to give it for free! Go beg the queen for free scraps, it's her fault the city is starving now! But leave our storming supplies alone!"
}

Allow \npcTroublemaker{} to mutter some awkward apologies and attract some shamespren. Unless the PCs act, proceed to the next interaction.



\subsection{\npcGrandma{}'s Demand}

Before you finish the scene in \npcGrandma{}'s home, read the following:

\quoteBox{
	As you are gathering on the ground floor, you hear \npcGrandma{} shouting to \npcTroublemaker{}: \\ \\
	"You must get me a proper food this time, not this tallew filled with cremlings! You can eat that yourself! Go to the market next to the People's Hall, they know who I am! You hear me? Tomorrow! First thing in the morning!"
}

As the characters leave the house, \npcTroublemaker{} will see them off:

\quoteBox{
	"Bye! And once again, thank you for today! Sorry for \npcGrandma{}, she is just... well... she's just \npcGrandma{}... But I know she liked you! Feel free to visit us anytime you need! \\ \\
	"By the way, will you join me tomorrow morning at the market? The one next to the People's Hall! Please?"
}

Once the characters are invited to witness the events of \fullref{sec:peoplesHall}, you can end this section.



\subsection{Aftermath}

Once the players complete their interactions in this scene, you can proceed to \fullref{sec:peoplesHall}.

